\renewcommand{\footerfilename}{Postgresql.tex}

\secton{PostgreSQL}

Documentation here is based on Postgresql 14.

PostgreSQL uses an OpenJDBC driver, available for download as
described at the PostgreSQL web site. This allows TOPCAT to load/save
tables directly from/to PostgreSQl. \index{TOPCAT!PostgreSQL}
\index{PostgreSQL!TOPCAT}.  TOPCAT, ds9, aladin, glueviz and if the
wind is favorable, the Chorimum browser, will allow the broadcasting
of data between themselves using SAMP.
\index{SAMP!TOPCAT!PostgreSQL}.

\subsection{Data from TOPCAT TAP Sources}

There are over 26,000 databases in VizieR. Most are not
relative, dated, etc -- but it is possible to grab
data related to one or more observing fields from places
like the GAIA EDR3, etc. 

This breaks down into two sub-problems:

\vspace{-.15cm}
\begin{enumerate}\addtolength{\itemsep}{-0.5\baselineskip}
%\setcounter{enumi}{N}
   \item   Get the data from wherever into TOPCAT
   \item   Store the data from TOPCAT into PostgreSQL.
\end{enumerate}


Data within TOPCAT is case sensitive, for example the Johnson R filter
and the Sloans filter 'r' are different.  These need to be
disambiguated when the case-sensitivity is removed. Lately TOPCAT is
add something like a \_x[a..z]* to do this, or some other means.
So R becomes R\_x and r becomes r\_xa. 

You can use the ALTER statement to achieve your naming desires.

A few cases pop up where the entered column names align with
PostgreSQL keywords, DEC is a good example. You can use
something like DEC\_SEXA or DEC\_d to to the trick. Taylor
the changes to match your program's requirements.

It is possible, within reason, to cross-match your list of
things by RA and Dec using the TOPCAT cross-match (with simbad)
feature. This brings up tonnes of columns. 

TOPCAT has the ability to directly query thousands of astronomical
databases. There are several paths to this, the best is the ``Select
From'' icon. The service and tables there in are excellent.






Some care has to be made to avoid the column name 'dec' for Declination, as
this is a SQL keyword. Case is not preserved making galactic longitude 'b'
and the Johnson filter 'B' collide in name space. It is not uncommon
to see BMag as a header. Some degree of conformance between tables is to
be expected. SIMBAD calls things oname, otype, and it seems natural to
extend this to include ora and odec as aliases for RA and Dec.


/etc/postgresq/14/main

/var/log/postgresql    postgresql-14-main.log



TOPCAT open /usr2/home/wayne/Observations/Data/wrstars.csv


\subsection{Administration}


TOPCAT reports error with standard installation of PostgreSQL


When TOPCAT saves a table, it does not add a sequence number. The
sequence number will precisely identify a row to correct any
mistakes etc. The best way is to save the TOPCAT table as a tmp
name, then use a ``CREATE...SELECT'' statement to create the
table with the desired name, add the sequence and prepare for
life in the future. Should the table already have the name you
want; you can simply rename the table. E.g.: TOPCAT save a
table called rawfoo, you want foo: Create table foo as select * from
rawfoo; Drop table if exists rawfoo; IF the table was saved as foo,
the ALTER TABLE foo RENAME TO rawfoo; and proceed as above.

\begingroup \fontsize{10pt}{10pt}
\selectfont
\begin{verbatim} 
/*****************************************************************************
* wr_all
*****************************************************************************/
DROP TABLE    IF EXISTS wr_all;
DROP SEQUENCE IF EXISTS wr_all_sequence;
CREATE SEQUENCE         wr_all_sequence START 100000;

CREATE TABLE wr_all as select * from wrall
COMMENT ON TABLE wr_all is 'Wolf Rayet Catalog';

\end{verbatim}
\endgroup



\subsection{Other PostgreSQL Tricks}

\begingroup \fontsize{10pt}{10pt}
\selectfont
\begin{verbatim} 
-- For a field named refstar, that has a ' n' at the end of the field:
-- remove the ' n'
update wrall_23dec set refstar = regexp_replace(refstar,' n','') where refstar like '% n';
select * from wrall_23dec where refstar != '';


select top 100 t1.main_id as "WR_ID", basic.main_id as "simbadname", 
    alltypes.otypes, ids.ids
from basic
join alltypes      on alltypes.oidref = basic.oid
join ids           on ids.oidref      = basic.oid
join TAP_UPLOAD.t1 on t1.main_id      = basic.main_id

-- Rename the table from some query blah to wr_types_aliases_raw
-- TOPCAT upload wr_types_aliases_raw to postgres

DROP TABLE    IF EXISTS wr_types_aliases;
DROP SEQUENCE IF EXISTS wr_types_aliases_sequence;
CREATE SEQUENCE         wr_types_aliases_sequence START 100000;

CREATE TABLE wr_types_aliases as select * from wr_types_aliases_raw
COMMENT ON TABLE wr_types_aliases_raw is 'Wolf Rayet Catalog alias names and types';


-- ------------------------------------------------------------------
-- Copy the valuable wrall_23dec table data as wr_new !!!
-- ------------------------------------------------------------------


DROP TABLE    IF EXISTS wr_new;
DROP SEQUENCE IF EXISTS wr_new_sequence;
CREATE SEQUENCE         wr_new_sequence START 100000;

CREATE TABLE wr_new as select * from wrall_23dec;
COMMENT ON TABLE wr_new is 'Wolf Rayet Catalog alias names and types';



alter table wr_new add column aliases  text[];
alter table wr_new add column types    text[];

-- look at incomming raw data
select * from wr_types_aliases_raw limit 10;

-- try the function on for size
select regexp_split_to_array(otypes,'|','e') from  wr_types_aliases_raw limit 10;
 
-- update one field at a time
UPDATE wr_new
   SET aliases = regexp_split_to_array(b.ids,'|','e') 
FROM wr_new as w
LEFT JOIN wr_types_aliases_raw AS b ON w.main_id = b.wr_id
;

UPDATE wr_new
   SET types   = regexp_split_to_array(b.otypes,'|','e')
FROM wr_new as w
LEFT JOIN wr_types_aliases_raw AS b ON w.main_id = b.wr_id
;


select * from wr_new limit 1;


select a.main_id, b.wr_id from wrall_23dec as a
   join wr_types_aliases_raw as b on b.wr_id = a.main_id;


select regexp_match(array_to_string(aliases,'|'),'.*(Gaia DR3 [ 0-9]+).*') as 'Gaia' from wr_new limit 1;
-- returns  {"Gaia DR3 524101256981955200"}

from astroquery.mast import Observations
obs_table = Observations.query_criteria(dataproduct_type=["spectrum"], obs_collection='HST', objectname='HD  93162')
print(obs_table[0:4])
if(o['intentType'] == 'science' and o['provenance_name'] == 'CALSTIS'):
   print(o)

data_products_by_obs = Observations.get_product_list(obs_table[1:3])

# Download only the science products
manifest = Observations.download_products(data_products_by_obs, productType="SCIENCE")
Downloading URL https://mast.stsci.edu/api/v0.1/Download/file?uri=mast:HST/product/z2u91502t_d0f.fits to ./mastDownload/HST/z2u91502t/z2u91502t_d0f.fits ...
|================================================================================================================|  57k/ 57k (100.00%)         0s
Downloading URL https://mast.stsci.edu/api/v0.1/Download/file?uri=mast:HST/product/z2za0102t_d1f.fits to ./mastDownload/HST/z2za0102t/z2za0102t_d1f.fits ...
|================================================================================================================|  14k/ 14k (100.00%)         0s


\end{verbatim}
\endgroup

\subsection{Schemas}





PostgreSQL uses \dhl{schema} to separate collections of tables within
a database. Databases maintain a hard and fast separation --
permitting several users to use the same engine with no knowledge or
access to the others' databases.

To view schemas within a database:

\begingroup \fontsize{10pt}{10pt}
\selectfont
%%\begin{Verbatim} [commandchars=\\\{\}]
\begin{verbatim} 
\dn

show search_path;
set search_path to chile;

\end{verbatim}
\endgroup
%% \end{Verbatim}

A schema is created with:

